\documentclass[11pt]{article}

\usepackage[a4paper,margin=2.5cm]{geometry}
\usepackage{amsmath}
\usepackage{amssymb}

\title{Construction Explained}
\author{Kirill Shmakov}
\date{\today}

\begin{document}

\maketitle

\section{Goal}

The idea of this project is to augment the algorithm presented
in~\cite{rokach2013mlcircuit} with an ML-backed model for complexity estimation
of decision trees, used for the heuristic choice of element.

\section{Algorithm}

Initially we have a boolean function $g$ with $n$ primary inputs $\bar{x} = \left(x_1, x_2, \dots, x_n\right)$.

Our goal is to construct a circuit representation of this function using logical
elements from $\mathcal{K} = \left\{ K_1, K_2, \dots, K_m \right\}$.

We use a greedy heuristic algorithm to find some implementation. This algorithm operates
on a pair $(S, C)$, where $S$ is the set of slots and $C$ is the current unfinished
circuit.

Initially $S_0 = \left\{x_1, x_2, \dots, x_n\right\}$ and $C_0$ is the empty
circuit.

At the iteration step for the current pair $(S, C)$ the algorithm iterates over the
elements $K$ of $\mathcal{K}$ and attaches each of them to the current slots $S$ in
all possible ways. The output of the element produces a new slot $y$. For every such
attachment we estimate a complexity score of our function $\text{Score}\left[g \mid \bar{x}, y\right]$.
The attachment that gives the minimal value of the score function is selected, and we
update our pair $(S, C)$.

The score function is the size of the minimal decision tree needed to compute the
function $g$ from the value table of $\bar{x}, y$: we pick the $y$ that maximally
helps us to reduce the remaining complexity.

\section{Model}

For $n \leqslant 12$ we can compute $\text{Score}\left[g \mid \bar{x}, y\right]$ as the size
of the minimal decision tree (the truth table is not very big, so dynamic programming
is applicable).

For larger $n$ this quickly becomes impossible, and we need to use an ML model $S_1$
for approximate computation.

$S_1$ accepts two boolean functions $g, f : \{0,1\}^n \to \{0,1\}$ as its arguments;
we write $S^n_1$ when we need to track the arity of the input functions.
It outputs a pair $(d, s)$ of real numbers, $d, s \in [0, n]$. Here $d$ is the depth
of the minimal-depth decision tree, and $s$ is $\log_2 (\text{IN} + 1)$, where
$\text{IN}$ is the number of internal nodes in the minimal-size decision tree; these
two trees need not be the same.

We also need another model $S_2$, which takes a boolean function $g : \{0,1\}^n \to \{0,1\}$
and a subset $X \subseteq \{0,1\}^n$ as its arguments. Its outputs $(d, s)$ have the
same semantics, for the decision trees needed to compute $g$ restricted to $X$.

For training $S_1$ and $S_2$ we use exact values when possible. Otherwise we compute
the target value as follows.

\subsection{Targets for $S_1$}

We select the target for $S^n_1\left[g \mid f\right]$ to be the minimum over the
following reductions:
\begin{itemize}
    \item through each primary input $x_i$:
        $$
            S^{n-1}_1 \left[g \downarrow_{x_i = 0} \mid f \downarrow_{x_i = 0} \right], \qquad
            S^{n-1}_1 \left[g \downarrow_{x_i = 1} \mid f \downarrow_{x_i = 1} \right]
        $$
    \item through $f$:
        $$
            S^n_2 \left[g \mid f^{-1}(0) \right], \qquad
            S^n_2 \left[g \mid f^{-1}(1) \right]
        $$
\end{itemize}

Here $h \downarrow_{x_i = \alpha}$ for $\alpha \in \{0, 1\}$ denotes the restriction of
$h$ obtained by fixing $x_i = \alpha$, viewed as a function of the remaining $n - 1$
inputs.

\subsection{Targets for $S_2$}

We select the target for $S^n_2\left[g \mid X\right]$ to be the minimum over the
reductions through each primary input $x_i$:
$$
    S^{n-1}_2 \left[g \downarrow_{x_i = 0} \mid X \downarrow_{x_i = 0} \right], \qquad
    S^{n-1}_2 \left[g \downarrow_{x_i = 1} \mid X \downarrow_{x_i = 1} \right]
$$
where $X \downarrow_{x_i = \alpha}$ is the projection of $X \cap \{x_i = \alpha\}$ onto
the coordinates other than $x_i$.

\bibliographystyle{plain}
\bibliography{references}

\end{document}
